\documentclass[11pt]{article}


\setlength{\parindent}{0pt}
\usepackage{xltxtra}
\usepackage{hyperref}
\hypersetup{hidelinks}
\usepackage{url}
\urlstyle{tt}
\usepackage{xcolor}
\definecolor{CVBlue}{RGB}{23,110,191}
\usepackage{calc}
\usepackage{graphicx}
\usepackage{tikz}
\usetikzlibrary{calc}
\usepackage{fontspec}
\usepackage{xeCJK}
\usepackage{enumitem}
\CJKsetecglue{} %% 取消中文与数字之间的间隙


%% 主文档字体设置
\setmainfont[
    Path = fonts/Main/,
    Extension = .otf,
    BoldFont = texgyretermes-bold.otf, % 加粗字体
]{texgyretermes-regular.otf} % 正文字体

% 中文字体设置
\setCJKmainfont[
    Path = fonts/hansans/,
    Extension = .ttf,
    BoldFont = NotoSansSC-Bold.ttf, % 加粗字体
]{NotoSansSC-Regular.otf} % 正文字体


\usepackage{relsize}
\usepackage{xspace}

% 使用 fontawesome（部分图标）
\usepackage{fontawesome} 

% A4纸，上下左右边距
\usepackage[
    a4paper,
    left=1.2cm,
    right=1.2cm,
    top=1.5cm,
    bottom=1cm,
    nohead
]{geometry}

\renewcommand{\baselinestretch}{1.5} % 行间距设为1.5

\usepackage{titlesec}
\usepackage{enumitem}
\setlist{noitemsep} % 取消列表项间的额外间距
%\setlist{nosep} % 取消所有垂直间距
\setlist[itemize]{topsep=0.25em, leftmargin=*}
\setlist[enumerate]{topsep=0.25em, leftmargin=*}

% --- 用于控制【不同项目之间】的垂直距离 ---
\newlength{\interProjectSpacing}
\setlength{\interProjectSpacing}{0.9em} % <--- 在此调整项目之间的距离
\newcommand{\projectsep}{\vspace{\interProjectSpacing}}

% --- 用于控制【项目标题】与下方【项目描述】的距离 ---
\newlength{\intraProjectTitleSep}
\setlength{\intraProjectTitleSep}{0.4em} % <--- 在此调整标题和描述的距离
\newcommand{\titlebreak}{\\[\intraProjectTitleSep]}

% --- 用于控制【项目描述】与下方【要点列表】的距离 ---
\newlength{\intraProjectListTopSep}
\setlength{\intraProjectListTopSep}{0.2em} % <--- 在此调整描述和列表的距离

% =======================================================================


\titleformat{\section}         % 定制 \section 命令 
{\large\bfseries\raggedright} % 将 section 标题设置为大号、粗体且左对齐
{}{0em}                      % 可用于为所有 section 添加前缀（如“章节...”）
{}                           % 可用于在标题前插入代码
[{\color{CVBlue}\titlerule}]  % 在标题后插入一条横线
\titlespacing*{\section}{0cm}{*1.6}{*1.2}



\begin{document}
\pagenumbering{gobble}

%%%% 利用tikz来定位照片
\begin{tikzpicture}[remember picture, overlay] 
    \node[anchor = north east] at ($(current page.north east)+(-2cm,-1.2cm)$) {\includegraphics[height=3cm]{avatar.jpg}};
  \end{tikzpicture}%
  %%%% 利用tikz来定位学校Logo，这里只在第一页显示
  \begin{tikzpicture}[remember picture, overlay] 
    \node[anchor = north west] at ($(current page.north west)+(0.5cm,+1.0cm)$) {\includegraphics[height=6cm]{zju.png}};
  \end{tikzpicture}%
\centerline{\LARGE\bfseries{刘思琪}} 

\centerline{\normalsize{\faPhone\ 158-1731-3220 \quad \faEnvelopeO\ \href{mailto:3230105088@zju.edu.cn}{3230100000@zju.edu.cn}}} 
    
\section{\makebox[\widthof{\faGraduationCap}][c]{\color{CVBlue}\faGraduationCap}\ 教育背景}    
\textbf{浙江大学} \hfill 2023.9 -- 至今\\[0.5em] % 标题和正文间加一点距离
自动化（控制）\quad 大三 
\begin{itemize}[nosep]
    \item 相关课程：《自动控制原理》、《现代控制理论》、《空中机器人》、《人工智能与机器学习》、《运动控制》
\end{itemize}

\section{\makebox[\widthof{\faUsers}][c]{\color{CVBlue}\faUsers}\ 项目经历}

% --- 第一个项目 ---
% 将标题行末尾的 \\ 替换为 \titlebreak 命令
\textbf{基于自然语言引导的视觉语义导航} \hfill 2025.06 -- 至今 \titlebreak
项目描述：为面向移动机器人的自然语言视觉语义导航系统，融合 ROS、YOLO 与视觉语言模型，实现从自然语言指令到自主导航执行的多模态闭环。
% 在 itemize 的选项中，使用 topsep=\intraProjectListTopSep 来控制上边距
\begin{itemize}[nosep, topsep=\intraProjectListTopSep]
    \item \textbf{系统与平台}：基于 \textbf{Ubuntu + ROS1 Noetic + Gazebo} 搭建完整机器人仿真平台，完成 \textbf{TurtleBot3} 模型部署、传感器配置、SLAM 建图与导航栈集成，构建可复现的语义导航实验环境。
    \item \textbf{多模态感知与理解}：集成 \textbf{YOLO} 实现环境中目标的实时检测与位姿输出，并引入 \textbf{视觉语言模型（VLM-Qwen）}，实现图像描述、语义问答与多模态推理，提升系统整体语义理解能力。
    \item \textbf{语义导航与规划}：设计基于 \textbf{LLM} 的自然语言指令解析与动作规划框架，实现从\textbf{自然语言输入 → 语义理解 → 路径规划 → 导航执行}的完整闭环，支持“找到沙发 / 冰箱”等语义导航任务。
\end{itemize}


% 使用 \projectsep 命令来分隔两个项目
\projectsep

% --- 第二个项目 ---
\textbf{基于数学优化的超市布局利润最大化研究 (Agent)} \hfill 2025.06 \titlebreak
项目描述：基于整数规划与动态规划构建超市商品布局与运营优化模型，通过考虑人流动态与连通性约束，实现长期利润最大化决策。
\begin{itemize}[nosep, topsep=\intraProjectListTopSep]
    \item \textbf{问题建模}：将超市抽象为 \textbf{5×5 网格布局模型}，引入初始人流量矩阵、商品购买率、单价及人流影响系数等参数，建立以\textbf{利润最大化}为目标的数学优化模型。
    \item \textbf{布局优化方法}：设计两种商品连通性约束方案——\textbf{基于图论的严格连通性约束}与\textbf{基于最小外接矩形的罚函数约束}，在保证合理布局的同时提升模型可解性与实际适用性。
    \item \textbf{人流动态建模}：考虑商品布局对顾客行为的反馈效应，构建\textbf{人流量动态更新机制}，重新定义全局人流矩阵，使模型更贴近真实超市运营场景。
    \item \textbf{长期规划与求解}：结合 \textbf{整数规划、非线性规划} 与 \textbf{动态规划} 方法，制定 \textbf{3 年超市装修与维护策略}，并使用 \textbf{Gurobi} 求解器获得最优利润方案。
\end{itemize}

% 使用 \projectsep 命令来分隔两个项目
\projectsep

\textbf{工控信息物理系统（ICPS）数据安全防护方案设计} \hfill 2025.06 \titlebreak
项目描述：构建面向 ICPS 的分层安全通信架构，通过标准密码算法组合实现多层级数据机密性、完整性与身份认证保障。
\begin{itemize}[nosep, topsep=\intraProjectListTopSep]
    \item \textbf{系统建模与需求分析}：基于典型 \textbf{工业控制信息物理系统（ICPS）分层架构}，分析管理层、监控层与控制层的通信特性与安全威胁，明确不同层级对\textbf{机密性、完整性与身份认证}的差异化需求。
    \item \textbf{分层加密方案设计}：遵循\textbf{分层防护、按需加密}原则，设计纵向与横向相结合的安全通信方案，组合采用 \textbf{TLS 1.3、AES-256-GCM 与 HMAC-SHA256} 实现多层级数据防护。
    \item \textbf{安全机制实现}：基于 \textbf{Python} 实现 TLS 加密通信、对称加密与消息认证机制，完成证书管理、加密传输、完整性校验等核心功能的原型开发与集成。
\end{itemize}


\section{\makebox[\widthof{\faCogs}][c]{\color{CVBlue}\faCogs}\ 技术栈}
\begin{itemize}[nosep]
    \item \textbf{编程语言：} \textbf{Go}, Python, C++
    \item \textbf{开发工具：}  Git, MakeFile,LaTex
    \item \textbf{操作系统：} Linux
\end{itemize}
    
\section{\makebox[\widthof{\faInfo}][c]{\color{CVBlue}\faInfo}\ 其他}
\begin{itemize}[parsep=0.5ex]
    \item \textbf{兴趣爱好：} 运动、看电影、旅行
    \item \textbf{英语水平：} CET-4, CET-6,IELTS
\end{itemize}
\end{document}
